\documentclass[12pt,a4paper,oneside]{article}
\usepackage[brazil]{babel}
\usepackage[utf8]{inputenc}
\usepackage{olymp}

\setcounter{problem}{1}

\begin{document}

\begin{problem}{Saldo de Gols}{saldo}{2}

O saldo de gols é um dos critérios de desempate usados em tabelas de classificação de futebol. O saldo de gols é a diferença entre o número de gols marcados e o número de gols sofridos. Por exemplo, se um time fez um total de 4 gols nas partidas já disputadas e não sofreu nenhum, seu saldo de gols é $4-0 = 4$. Se um time fez 1 gol e sofreu 3, sem saldo de gols é $1-3 = -2$.

Faça um programa para ler os resultados das partidas disputadas por um time e informar seu saldo de gols.

%\Notes
%
%Observações, se tiver.

\InputFile

A entrada começa com um inteiro $N$, o número de partidas já disputadas por um time. As $N$ linhas seguintes contém 2 inteiros cada, $M_i$ e $S_i$, respectivamente o número de gols marcados e sofridos pelo time na $i$-ésima partida. Restrições: $1 \leq N \leq 38$; $0 \leq M_i, S_i \leq 20$.


\OutputFile

Seu programa deve escrever uma linha informando o saldo de gols do time.

% \Notes
% observações

\Example

\begin{example}
\exmp{
3
0 1
1 2
1 2
}{
-3
}%
\end{example}

\begin{example}
\exmp{
3
0 0
4 0
1 1
}{
4
}%
\end{example}

\begin{example}
\exmp{
6
5 1
1 0
1 1
0 2
1 2
0 1
}{
1
}%
\end{example}


\textit{Os dois primeiros exemplos são resultados de times mineiros na Copa Sul-Americana 2025 até a data que a questão foi dada em aula; no último, resultados do Brasil nas eliminatórias da Copa.}

\end{problem}

\end{document}
